\section{System Model}

\subsection{Scope}

The model covers a single asset and limit orders only. Each order carries an
identifier, account, side, limit price, quantity, and timestamp. The minimum
version includes only submission and matching transitions; cancellation and
partial fills are treated as stretch goals because they substantially increase
the amount of residual-state reasoning required.

\subsection{State}

The engine state consists of:

\begin{itemize}[leftmargin=1.5em]
\item a buy-side order book;
\item a sell-side order book;
\item a trade log;
\item per-account cash and asset holdings; and
\item any auxiliary bookkeeping needed for timestamps or order identity.
\end{itemize}

For proof convenience, each side of the book can be represented either as a
sorted sequence or as a relational ranking structure. The current plan
prioritises a sorted representation because it makes price-time priority
visible in the state and aligns naturally with helper lemmas about insertion
and deletion.

\subsection{Transitions}

The abstract transition system contains at least the following steps:

\begin{itemize}[leftmargin=1.5em]
\item \texttt{submit}: add a new limit order to the appropriate side of the
book while preserving book well-formedness;
\item \texttt{match}: execute one trade between the current best bid and best
ask when they are compatible; and
\item optionally \texttt{cancel}: remove a live order from the book.
\end{itemize}

The paper focuses on the first two transitions. If cancellation is included,
it will be added only after the base safety proof is stable.
